% ============================================================================
%  Observational Uniqueness for Zero-Parameter Frameworks
%  Jonathan Washburn, March 2026
% ============================================================================
\documentclass[11pt]{article}

\usepackage[margin=1in]{geometry}
\usepackage{microtype}
\usepackage{lmodern}
\usepackage[T1]{fontenc}
\usepackage[utf8]{inputenc}
\usepackage{amsmath, amssymb, amsthm, mathtools}
\usepackage[colorlinks=true,linkcolor=blue,citecolor=teal,urlcolor=magenta]{hyperref}
\IfFileExists{cleveref.sty}{%
  \usepackage[nameinlink,capitalise]{cleveref}%
}{%
  \providecommand{\cref}[1]{\ref{##1}}%
  \providecommand{\Cref}[1]{\ref{##1}}%
}%

\newcommand{\R}{\mathbb{R}}
\newcommand{\N}{\mathbb{N}}
\newcommand{\Rp}{\mathbb{R}_{>0}}
\newcommand{\Jcost}{J_{\mathrm{cost}}}
\newcommand{\SQ}{\mathrm{SQ}}
\newcommand{\cZ}{\mathcal{Z}}
\newcommand{\Odim}{\mathcal{O}_{\mathrm{dim}}}

\newtheorem{theorem}{Theorem}[section]
\newtheorem{lemma}[theorem]{Lemma}
\newtheorem{proposition}[theorem]{Proposition}
\newtheorem{corollary}[theorem]{Corollary}
\theoremstyle{definition}
\newtheorem{definition}[theorem]{Definition}
\newtheorem{example}[theorem]{Example}
\newtheorem{openproblem}[theorem]{Open Problem}
\theoremstyle{remark}
\newtheorem{remark}[theorem]{Remark}

\title{Observational Uniqueness for Zero-Parameter Frameworks\\[4pt]
\large Quotient Collapse Under the Admissibility Posture,\\
and the Open Programme for Full Inevitability}
\author{Jonathan Washburn}
\date{March 2026}

\hypersetup{
  pdftitle={Observational Uniqueness for Zero-Parameter Frameworks},
  pdfauthor={Jonathan Washburn},
}

\begin{document}
\maketitle

% ===================================================================
\begin{abstract}
We prove that any physics framework satisfying five structural
constraints---inhabitedness, zero adjustable parameters, self-similarity,
an admissible cost functional obeying a coherent composition law, and
observables extracted via cost minimization---is forced into a unique
observational state.
Formally: the quotient of the state space by observational
indistinguishability collapses to a singleton.
The preferred scale is forced to the golden ratio
$\varphi = (1+\sqrt{5})/2$, and the cost functional is uniquely
$\Jcost(x) = \tfrac{1}{2}(x + x^{-1}) - 1$.

Three of the five assumptions (self-similarity, the composition law, and
cost-minimum semantics) are \emph{adopted} axioms whose derivation from
the zero-parameter posture alone is an open programme.
We name six structural bridges that would close this programme,
characterize why each is hard, and identify the derivation of the
composition law from ledger structure as the deepest open problem.
The composition law is not a bespoke equation: we show it is equivalent
to the d'Alembert functional equation, classified by
Acz\'el (1966).

All proved results are formally verified in Lean~4 with zero
\texttt{sorry}.
Core results are conditional on one external axiom (the Acz\'el
classification theorem for continuous d'Alembert solutions).
\end{abstract}

\tableofcontents
\newpage

% ===================================================================
\section{Introduction}\label{sec:intro}
% ===================================================================

\subsection{The question}

The Standard Model of particle physics has its symmetry group
fixed by gauge consistency, but its 19 free parameters---coupling
constants, masses, mixing angles---must be measured, not derived.
What would happen if one demanded \emph{zero} adjustable
dimensionless constants?

This paper answers a precise version of that question.
We define a minimal interface for a ``physics framework''
(a state space, evolution, and measurement map),
formalize ``zero parameters'' as algorithmic describability
of the state space,
and prove:

\begin{quote}
\emph{Under five explicit structural assumptions, every admissible
framework has exactly one observational state on the quotient.}
\end{quote}

\noindent
The five assumptions divide into two tiers (\cref{sec:assumptions}):
two that are forced by the zero-parameter posture (Tier~1),
and three that are adopted as structural axioms whose derivation
from Tier~1 alone is an open programme (Tier~2).

\subsection{What this paper proves and what it does not}

\begin{center}
\renewcommand{\arraystretch}{1.2}
\begin{tabular}{|p{3.5cm}|p{4.5cm}|p{4.5cm}|}
\hline
\textbf{Level} & \textbf{Proved} & \textbf{Open} \\
\hline
\textbf{Tier 1} (A1)+(A2) &
  Discreteness, comparison structure, countable skeleton &
  --- \\
\hline
\textbf{Bridges B1--B4} &
  Fibonacci minimality among self-similar schemas &
  (HS)$\Rightarrow$B1; RCL derivation (B2, deepest);
  ratio interface B3; ground-state B4 \\
\hline
\textbf{Tier 2} (A1)--(A5) &
  $\sigma = \varphi$,\; $J = \Jcost$,\;
  $\mu_F$ uniform,\; $\SQ(F) \simeq \mathbf{1}$ &
  --- \\
\hline
\textbf{Reduction} &
  Both $F$ and canonical framework have trivial quotient &
  B5: prediction map;\; B6: value equality \\
\hline
\end{tabular}
\end{center}

\medskip\noindent
The core contribution is the ``Tier~2'' row.
Full inevitability requires closing bridges B1--B6
(\cref{sec:bridges}).

\subsection{Relation to existing work}

Wigner (1939) classified particles by Poincar\'e representations.
Coleman--Mandula (1967) showed the $S$-matrix symmetry group
is uniquely determined.
Both share the logic: \emph{structural constraints collapse
the space of admissible theories}.
Our result is of this type, applied to the zero-parameter posture
rather than to relativistic symmetry.

The cost-uniqueness theorem (T5) was proved independently by
Washburn--Zlatanovi\'c~\cite{WZ2026} and
Washburn--Rahnamai~Barghi~\cite{Barghi2026};
the d'Alembert reduction was formalized in~\cite{DAlembert2026}.
The classical reference for the d'Alembert functional equation is
Acz\'el~\cite{Aczel1966}.

\subsection{Roadmap}

\Cref{sec:framework} defines the framework interface and the
observational quotient.
\Cref{sec:assumptions} states the five assumptions and their tier
classification.
\Cref{sec:necessity} provides the necessity stack justifying
why the assumptions are not ad hoc.
\Cref{sec:dalembert} exposes the d'Alembert structure of the
composition law---the strongest argument that (A4) is forced.
\Cref{sec:t5} proves cost uniqueness ($J = \Jcost$).
\Cref{sec:main} proves the main theorem (quotient collapse).
\Cref{sec:bridges} names the six open bridges and
characterizes the open derivation programme.
\Cref{sec:falsifiers} lists concrete falsifiers.

% ===================================================================
\section{Frameworks and the Observational Quotient}\label{sec:framework}
% ===================================================================

\begin{definition}[Physics Framework]\label{def:framework}
A \emph{physics framework} is a tuple
$F = (S_F, T_F, O_F, \mu_F)$
where $S_F$ is a nonempty set (states),
$T_F : S_F \to S_F$ is an evolution operator,
$O_F$ is a set (observables),
and $\mu_F : S_F \to O_F$ is a measurement map.
\end{definition}

\begin{definition}[Algorithmic Specification]\label{def:alg-spec}
A set $S$ admits an \emph{algorithmic specification} if there exist
a finite binary string $d$,
an enumeration $g : \N \to \{0,1\}^* \cup \{\bot\}$,
and a decoder $\psi : \{0,1\}^* \to S \cup \{\bot\}$
such that for every $s \in S$ there exist $n \in \N$ and
$c \in \{0,1\}^*$ with $g(n) = c$ and $\psi(c) = s$.
\end{definition}

\begin{definition}[Zero-Parameter Framework]\label{def:zero-param}
$F$ has \emph{zero parameters} (written $F \in \cZ$) if
$S_F$ admits an algorithmic specification.
\end{definition}

\begin{definition}[Observational Equivalence]\label{def:obs-equiv}
For a framework $F$, define
$s_1 \sim_F s_2 \;\Longleftrightarrow\; \mu_F(s_1) = \mu_F(s_2)$.
\end{definition}

\begin{proposition}\label{prop:equiv-rel}
$\sim_F$ is an equivalence relation.
\end{proposition}
\begin{proof}
Reflexivity, symmetry, and transitivity of equality.
\end{proof}

\begin{definition}[Quotient State Space]\label{def:quotient}
$\SQ(F) := S_F / {\sim_F}$, the set of observational equivalence
classes.
We say the quotient \emph{collapses} when $\SQ(F)$ is a subsingleton
(at most one element).
\end{definition}

\begin{lemma}[Uniform Observables Imply Quotient Collapse]
\label{lem:uniform-collapse}
If $\mu_F(s_1) = \mu_F(s_2)$ for all $s_1, s_2 \in S_F$,
then $\SQ(F)$ is a subsingleton.
\end{lemma}
\begin{proof}
Any two classes $[a], [b]$ satisfy $\mu_F(a) = \mu_F(b)$,
hence $a \sim_F b$, hence $[a] = [b]$.
\end{proof}

% ===================================================================
\section{The Five Assumptions}\label{sec:assumptions}
% ===================================================================

\begin{definition}[Model-Independent Assumptions]\label{def:assumptions}
The \emph{admissibility bundle} for a framework $F$ consists of:
\begin{enumerate}
\item[\textup{(A1)}] $F$ is \emph{inhabited}: $S_F \neq \emptyset$.
\item[\textup{(A2)}] $F$ has \emph{zero parameters}: $F \in \cZ$.
\item[\textup{(A3)}] $F$ admits a \emph{self-similarity structure}
  (\cref{def:self-sim}).
\item[\textup{(A4)}] $F$ has an \emph{admissible cost functional} $J$
  satisfying the six conditions of \cref{def:admissible-cost}.
\item[\textup{(A5)}] $F$'s observables are \emph{extracted via cost}
  with all states at the cost minimum
  (\cref{def:obs-from-cost,def:measure-by-cost}).
\end{enumerate}
\end{definition}

\medskip
\begin{center}
\renewcommand{\arraystretch}{1.15}
\begin{tabular}{lll}
\hline
\textbf{Assumption} & \textbf{Status} & \textbf{Tier} \\
\hline
(A1) Inhabited & Forced (any non-empty framework) & 1 \\
(A2) Zero parameters & Definitional (class $\cZ$) & 1 \\
(A3) Self-similarity & Adopted; derivation from (A2) open & 2 \\
(A4) Cost + composition law & Adopted; derivation from ledger open & 2 \\
(A5) Observables from cost & Adopted; ratio interface + ground-state open & 2 \\
\hline
\end{tabular}
\end{center}

\medskip\noindent
Tier~1 assumptions are consequences of the zero-parameter posture.
Tier~2 assumptions import structural data whose derivation from
(A1)--(A2) alone is the \emph{open derivation programme}
(\cref{sec:bridges}).
All theorems in this paper are conditional on the full bundle
(A1)--(A5).

% ===================================================================
\section{The Necessity Stack}\label{sec:necessity}
% ===================================================================

Each Tier~2 assumption is justified---though not yet derived---by
a chain of structural necessities:
\[
\text{Observables}
\;\Rightarrow\; \text{Comparison}
\;\Rightarrow\; \text{Discrete}
\;\Rightarrow\; \text{Conservation}
\;\Rightarrow\; \varphi.
\]
We summarize each arrow.

\subsection{Observables require comparison}

\begin{proposition}\label{prop:obs-comparison}
Any framework that derives non-trivial observables in a
zero-parameter setting without external oracles must internally
realize a reflexive, symmetric comparison mechanism that
distinguishes states with different observable values.
\end{proposition}

\begin{proof}[Argument]
The measurement map $\mu_F$ is a function $S_F \to O_F$.
If $\mu_F$ is non-constant, there exist $s_1, s_2$ with
$\mu_F(s_1) \neq \mu_F(s_2)$.
In a closed system with no external oracle, the detection
of this inequality must be realized by an internal function
$\delta : S_F \times S_F \to \{0,1\}$ satisfying
$\mu_F(s_1) \neq \mu_F(s_2) \Rightarrow \delta(s_1, s_2) = 1$.
Reflexivity ($\delta(s,s) = 1$) and symmetry
($\delta(s_1, s_2) = \delta(s_2, s_1)$) follow from equality
being an equivalence relation.
\end{proof}

\subsection{Zero parameters force discreteness}

\begin{proposition}\label{prop:discrete}
Any $F \in \cZ$ admits a surjection from a countable set onto $S_F$.
\end{proposition}
\begin{proof}
The algorithmic specification provides an enumeration
$g : \N \to \{0,1\}^*$ and a decoder $\psi$ covering every
$s \in S_F$.
The composite $\psi \circ g$ is a surjection from a countable
index set.
\end{proof}

\subsection{Discrete conservation forces balanced accounting}

\begin{proposition}\label{prop:ledger}
Given a locally finite discrete event system with a conserved flow,
the conservation law is equivalent to a balanced double-entry
condition on events.
\end{proposition}
\begin{proof}[Sketch]
Conservation at each event $e$ requires
$\sum_{v \to e} f(v) = \sum_{e \to w} f(w)$.
This is precisely the balance condition of a double-entry
accounting system.
The discrete, locally finite structure ensures the sums are
finite and well-defined.
\end{proof}

\subsection{Self-similarity forces $\varphi$}

\begin{definition}[Self-Similarity Structure]\label{def:self-sim}
A framework admits a \emph{self-similarity structure} if there
exist a preferred scale $\sigma > 1$ and three successive levels
$\ell_0, \ell_1, \ell_2 > 0$ with
$\ell_1 = \sigma \ell_0$,\;
$\ell_2 = \sigma \ell_1$,\; and
$\ell_2 = \ell_1 + \ell_0$.
\end{definition}

\begin{theorem}[Self-Similarity Forces $\varphi$]\label{thm:phi}
If $F$ admits a self-similarity structure with scale $\sigma$, then
\[
\sigma^2 = \sigma + 1,
\qquad\text{hence}\qquad
\sigma = \varphi := \frac{1 + \sqrt{5}}{2}.
\]
\end{theorem}
\begin{proof}
From the definitions:
$\ell_2 = \sigma^2 \ell_0$ and
$\ell_2 = (\sigma + 1)\ell_0$.
Since $\ell_0 > 0$, cancel to get $\sigma^2 = \sigma + 1$.
The polynomial $x^2 - x - 1$ has roots
$(1 \pm \sqrt{5})/2$.
Since $\sigma > 1$, we have $\sigma = (1 + \sqrt{5})/2 = \varphi$.
\end{proof}

% ===================================================================
\section{The d'Alembert Structure of the Composition Law}
\label{sec:dalembert}
% ===================================================================

This section exposes the composition law (A4) as a classical
functional equation, not a bespoke recognition-specific identity.
This is the strongest argument that (A4) is not ad hoc.

\subsection{The Recognition Composition Law}

\begin{definition}[Admissible Cost Functional]\label{def:admissible-cost}
A cost functional $J : \Rp \to \R$ is \emph{admissible} if:
\begin{enumerate}
\item \emph{Reciprocal symmetry:}
  $J(x) = J(x^{-1})$ for all $x > 0$.
\item \emph{Unit normalization:} $J(1) = 0$.
\item \emph{Strict convexity:} $J$ is strictly convex on $\Rp$.
\item \emph{Calibration:}
  $\frac{d^2}{dt^2}(J \circ \exp)\big|_{t=0} = 1$.
\item \emph{Continuity:} $J$ is continuous on $\Rp$.
\item \emph{Composition law (RCL):} For all $x, y > 0$,
\begin{equation}\label{eq:rcl}
J(xy) + J(x/y) = 2J(x)J(y) + 2J(x) + 2J(y).
\end{equation}
\end{enumerate}
\end{definition}

\subsection{Reduction to the d'Alembert equation}

Define the reparametrizations
\begin{equation}\label{eq:GH}
G(t) := J(e^t), \qquad H(t) := G(t) + 1 = J(e^t) + 1.
\end{equation}

\begin{proposition}[RCL $\Leftrightarrow$ d'Alembert]
\label{prop:dalembert-equiv}
$J$ satisfies \eqref{eq:rcl} if and only if $H$ satisfies the
d'Alembert functional equation:
\begin{equation}\label{eq:dalembert}
H(s + t) + H(s - t) = 2\,H(s)\,H(t),
\qquad s, t \in \R.
\end{equation}
Moreover, the admissibility conditions translate as follows:
\begin{itemize}
\item $J(1) = 0 \;\Leftrightarrow\; H(0) = 1$.
\item $J(x) = J(x^{-1}) \;\Leftrightarrow\; H(-t) = H(t)$
  \textup{(}evenness\textup{)}.
\item Continuity of $J$ on $\Rp \;\Leftrightarrow\;$
  continuity of $H$ on $\R$.
\end{itemize}
\end{proposition}
\begin{proof}
Substitute $x = e^s$, $y = e^t$ into~\eqref{eq:rcl} and expand
$K = J + 1$:
\begin{align*}
J(e^{s+t}) + J(e^{s-t})
&= 2J(e^s)J(e^t) + 2J(e^s) + 2J(e^t) \\
G(s+t) + G(s-t)
&= 2G(s)G(t) + 2G(s) + 2G(t) \\
(H(s+t) - 1) + (H(s-t) - 1)
&= 2(H(s)-1)(H(t)-1) + 2(H(s)-1) + 2(H(t)-1) \\
H(s+t) + H(s-t) - 2
&= 2H(s)H(t) - 2 \\
H(s+t) + H(s-t)
&= 2H(s)H(t). \qedhere
\end{align*}
\end{proof}

\subsection{The Acz\'el classification}

The d'Alembert equation \eqref{eq:dalembert} is classical.
Jean le Rond d'Alembert studied it in the eighteenth century
in connection with vibrating strings.
The complete classification of continuous solutions was established
by Acz\'el~\cite{Aczel1966}:

\begin{theorem}[Acz\'el Classification {\cite[Ch.~3]{Aczel1966}}]
\label{thm:aczel}
The continuous solutions of
$H(s+t) + H(s-t) = 2H(s)H(t)$ with $H(0) = 1$ are exactly:
\begin{enumerate}
\item[\textup{(i)}] $H(t) = 1$ for all $t$
  \;\textup{(}constant\textup{)},
\item[\textup{(ii)}] $H(t) = \cosh(\lambda t)$ for some
  $\lambda \neq 0$,
\item[\textup{(iii)}] $H(t) = \cos(\lambda t)$ for some
  $\lambda \neq 0$.
\end{enumerate}
\end{theorem}

\begin{remark}[Why the RCL is not ad hoc]\label{rem:not-adhoc}
The composition law~\eqref{eq:rcl} is not a bespoke equation
engineered to produce a desired answer.
It is the d'Alembert functional equation in multiplicative
coordinates---the unique symmetric composition law on
$(\Rp, \times)$ compatible with non-degeneracy.
The analogy with the parallelogram law is precise:
Jordan--von~Neumann (1935) proved that the parallelogram identity
uniquely characterizes inner-product spaces among normed spaces;
the d'Alembert identity uniquely characterizes $\cosh$ among
symmetric composition functions on abelian groups.
\end{remark}

\subsection{Why additive composition fails}

One might ask: why not $J(xy) = J(x) + J(y)$ (logarithmic cost)?
Logarithms satisfy $J(x) = J(x^{-1})$ only if $J \equiv 0$,
violating strict convexity.
The d'Alembert form arises specifically from combining reciprocal
symmetry with the requirement that the composition is quadratic
(not additive) in $J$.
Establishing \emph{why} the composition must be quadratic rather
than additive is Bridge B2 (\cref{sec:bridge-b2}).

% ===================================================================
\section{Cost Uniqueness}\label{sec:t5}
% ===================================================================

\begin{theorem}[T5: Cost Uniqueness]\label{thm:t5}
Let $J : \Rp \to \R$ be an admissible cost functional
(\cref{def:admissible-cost}).
Then for all $x > 0$:
\begin{equation}\label{eq:jcost}
J(x) = \Jcost(x) := \tfrac{1}{2}(x + x^{-1}) - 1.
\end{equation}
\end{theorem}

\begin{proof}
By \cref{prop:dalembert-equiv}, $H(t) = J(e^t) + 1$ is a
continuous, even solution of the d'Alembert equation with
$H(0) = 1$.

\emph{Step 1: Apply the Acz\'el classification
(\cref{thm:aczel}).}
The three cases are: (i)~$H = 1$, (ii)~$H = \cosh(\lambda t)$,
(iii)~$H = \cos(\lambda t)$.

\emph{Step 2: Eliminate cases (i) and (iii).}
Case~(i) gives $J \equiv 0$, contradicting strict convexity
(which requires a unique minimum, not a flat zero).
Case~(iii) gives $J(e^t) = \cos(\lambda t) - 1 \leq 0$
with a maximum at $t = 0$;
strict convexity requires a \emph{minimum} at $t = 0$, ruling
this out.

\emph{Step 3: Calibrate case (ii).}
$H(t) = \cosh(\lambda t)$ gives
$\frac{d^2}{dt^2} H\big|_{t=0} = \lambda^2$.
The calibration condition forces $\lambda^2 = 1$,
hence $\lambda = 1$ (taking $\lambda > 0$).

\emph{Step 4: Conclude.}
$J(e^t) = \cosh(t) - 1 = \frac{1}{2}(e^t + e^{-t}) - 1$.
Substituting $x = e^t$:
$J(x) = \frac{1}{2}(x + x^{-1}) - 1 = \Jcost(x)$.
\end{proof}

\begin{corollary}\label{cor:unique-min}
$\Jcost(x) = 0$ if and only if $x = 1$.
$\Jcost$ is strictly convex on $\Rp$ with unique minimum at $x = 1$.
\end{corollary}
\begin{proof}
$\Jcost(x) = \frac{1}{2}(x + x^{-1}) - 1 \geq 0$ by AM--GM,
with equality iff $x = x^{-1}$, i.e., $x = 1$ (since $x > 0$).
\end{proof}

% ===================================================================
\section{Observables from Cost and Derived Uniformity}
\label{sec:obs-cost}
% ===================================================================

\begin{definition}[Measure Determined by Cost]\label{def:measure-by-cost}
$F$ has \emph{measure determined by cost} $J$ if there exists
$r : S_F \to \Rp$ (the ratio extraction) such that
\[
r(s_1) = r(s_2) \;\Longrightarrow\; \mu_F(s_1) = \mu_F(s_2).
\]
\end{definition}

\begin{definition}[Observables from Cost]\label{def:obs-from-cost}
$F$ has \emph{observables from cost} if it has measure determined
by cost and every state lies at the cost minimum:
$J(r(s)) = 0$ for all $s \in S_F$.
\end{definition}

\begin{theorem}[Derived Uniformity]\label{thm:derived-uniform}
Let $F$ be a framework with an admissible cost functional $J$,
measure determined by cost via $r : S_F \to \Rp$,
and $J(r(s)) = 0$ for all $s$.
Then:
\begin{enumerate}
\item $J = \Jcost$ \;\textup{(}by T5\textup{)}.
\item $r(s) = 1$ for all $s$
  \;\textup{(}by \cref{cor:unique-min}\textup{)}.
\item $\mu_F$ is uniform: $\mu_F(s_1) = \mu_F(s_2)$
  for all $s_1, s_2$.
\end{enumerate}
\end{theorem}
\begin{proof}
(1)~\Cref{thm:t5}.
(2)~$\Jcost(r(s)) = 0$ and $\Jcost(x) = 0 \Leftrightarrow x = 1$
(\cref{cor:unique-min}).
(3)~$r(s_1) = 1 = r(s_2)$, so measure-determined-by-cost gives
$\mu_F(s_1) = \mu_F(s_2)$.
\end{proof}

% ===================================================================
\section{Main Theorem}\label{sec:main}
% ===================================================================

\begin{theorem}[Observational Uniqueness on the Quotient]
\label{thm:main}
Let $F$ be a physics framework satisfying
\textup{(A1)--(A5)} of \cref{def:assumptions}.
Then:
\begin{enumerate}
\item The preferred scale is $\varphi = (1 + \sqrt{5})/2$.
\item The cost functional is $J = \Jcost$.
\item The measurement map is uniform.
\item The quotient state space is a subsingleton:
  $\SQ(F) \simeq \mathbf{1}$.
\end{enumerate}
\end{theorem}

\begin{proof}
\emph{(1)}
By (A3) and \cref{thm:phi}, $\sigma^2 = \sigma + 1$ with
$\sigma > 1$ gives $\sigma = \varphi$.

\emph{(2)}
By (A4) and \cref{thm:t5}, the admissible cost functional
equals $\Jcost$ on $\Rp$.

\emph{(3)}
By (A5) and \cref{thm:derived-uniform}, cost-minimum semantics
and cost uniqueness force $r(s) = 1$ for all $s$, hence
$\mu_F$ is uniform.

\emph{(4)}
By \cref{lem:uniform-collapse}, uniform observables imply
$\SQ(F)$ is a subsingleton.
Since $F$ is inhabited (A1), $\SQ(F)$ is nonempty.
A nonempty subsingleton is equivalent to $\mathbf{1}$.
\end{proof}

\begin{corollary}[Framework Equivalence]\label{cor:equiv}
Let $\mathcal{C} = (\mathbf{1}, \mathrm{id}, \Odim, \mu_\mathcal{C})$
be the canonical one-state framework.
If $F$ satisfies \textup{(A1)--(A5)} and there exists an
observable-type equivalence $e : O_F \simeq \Odim$, then
\[
\SQ(F) \simeq \SQ(\mathcal{C}) \simeq \mathbf{1}
\qquad\text{and}\qquad
O_F \simeq O_\mathcal{C}.
\]
\end{corollary}

\begin{remark}[What the theorem does not claim]
The theorem does \emph{not} assert that the raw state space
$S_F$ is a singleton.
It asserts that all internal states are observationally
indistinguishable at the cost-extracted interface:
multiplicity of $S_F$ is gauge redundancy, not physical freedom.
\end{remark}

% ===================================================================
\section{The Open Derivation Programme: Six Bridges}
\label{sec:bridges}
% ===================================================================

The full inevitability claim---``any $F \in \cZ$ is forced into
the RS structure''---requires deriving (A3)--(A5) from
(A1)--(A2) alone.
This section names the six structural bridges.

\subsection{Bridge B1: self-similarity from zero parameters}
\label{sec:bridge-b1}

\Cref{thm:phi} proves: \emph{among} self-similar schemas,
$\varphi$ is forced.
It does not prove: a zero-parameter ledger \emph{must} be
self-similar.

\begin{openproblem}[Hierarchy Theorem]
Show that a discrete, zero-parameter ledger with at least two
levels of event composition and uniform inter-level scaling
necessarily satisfies a monic integer recurrence.
This, combined with \cref{thm:phi}, forces $\varphi$.
\end{openproblem}

The intermediate result ``(A2) + hierarchical structure
$\Rightarrow$ (A3)'' is well-posed.
Deriving hierarchical structure from (A2) alone remains open.

\subsection{Bridge B2: the RCL from ledger structure (deepest)}
\label{sec:bridge-b2}

The RCL~\eqref{eq:rcl} is adopted as an axiom in (A4).
\Cref{sec:dalembert} showed it is the d'Alembert equation---not
a bespoke formula.
The open problem is to \emph{derive} it:

\begin{openproblem}[Composition Rule Classification]
\label{op:rcl}
Among all continuous, symmetric composition rules
$f : \R \times \R \to \R$ satisfying
$J(xy) + J(x/y) = f(J(x), J(y))$ for a reciprocally symmetric
cost $J$ with $J(1) = 0$:
\begin{enumerate}
\item $f(0,0) = 0$ and $f(a, 0) = 2a$
  \;\textup{(}from setting $y = 1$\textup{)}.
\item If cost composition is associative
  \textup{(}derivable from the group structure of
  $(\Rp, \times)$\textup{)}, then
  $f(a,b) = 2(a+1)(b+1) - 2$,
  i.e., $J$ satisfies the RCL.
\end{enumerate}
Prove or disprove that associativity is forced by the
zero-parameter posture.
\end{openproblem}

This is the most significant open problem in the programme.

\subsection{Bridge B3: the ratio-extraction interface}
\label{sec:bridge-b3}

(A5) posits $r : S_F \to \Rp$.
This bundles three sub-claims:
\begin{enumerate}
\item \emph{Dimensionlessness:} all observables are ratios
  (no dimensionful constants in a zero-parameter framework).
\item \emph{One-dimensionality:} a single real per state
  (requires the ledger to have one independent conserved
  quantity).
\item \emph{Positivity:} the ratio lives on $\Rp$.
\end{enumerate}

\subsection{Bridge B4: the ground-state restriction}
\label{sec:bridge-b4}

The hypothesis $J(r(s)) = 0$ for all $s$ is not automatic.
By reciprocal symmetry, $r(S_F)$ is inversion-closed.
Under a \emph{structural determinacy} (SD) hypothesis---that
each ratio is uniquely determined by a computable,
$\mathbb{Z}/2\mathbb{Z}$-invariant predicate---any $x \neq 1$
requires at least 1 bit to select from $\{x, x^{-1}\}$,
contradicting zero parameters.
Deriving (SD) from (A2) alone is open.

\begin{example}[Concrete falsifier for automatic ground-state]
Let $S_F = \{s_0\}$, $r(s_0) = 2$.
Then $\Jcost(r(s_0)) = \Jcost(2) = \frac{1}{4} \neq 0$.
This one-state framework satisfies (A1), (A2), (A4), and
vacuously (A5-i), but violates the ground-state restriction.
\end{example}

\subsection{Bridge B5: the constructive prediction map}
\label{sec:bridge-b5}

\Cref{thm:main} proves $\mu_F$ is constant but does not identify
the constant.
A prediction map
$\mathcal{P} : (\Jcost, \varphi) \to \Odim$ would need to:
(i)~take $\Jcost$ and $\varphi$ as sole inputs,
(ii)~produce specific dimensionless quantities,
(iii)~have $O(1)$ description length, and
(iv)~be unique.
Existence (i)--(iii) can be established by formalizing the RS
ledger algorithm; uniqueness (iv) and value identification
remain open.

\subsection{Bridge B6: the observable-type equivalence}
\label{sec:bridge-b6}

\Cref{cor:equiv} requires $e : O_F \simeq \Odim$ as input.
After quotient collapse, $\mathrm{im}(\mu_F)$ is a singleton.
Any two singleton types are abstractly isomorphic,
so the type-level equivalence is trivially constructible.
Non-trivial content (value identification) reduces to B5.

% ===================================================================
\section{Falsifiers}\label{sec:falsifiers}
% ===================================================================

The theorem is a mathematical statement that cannot be refuted
within its hypotheses.
However, the \emph{inevitability framing} can be challenged by
showing that a hypothesis is too strong, too narrow, or physically
unwarranted.
Each of the following would constitute a substantive refutation:

\begin{enumerate}
\item \textbf{Alternative zero-parameter framework.}
  Construct $F \in \cZ$ that derives observables and does
  \emph{not} collapse on the quotient (non-uniform
  $\mu_F$ without free parameters).

\item \textbf{Non-$\varphi$ self-similarity.}
  Provide a parameter-free self-similar scaling whose preferred
  scale is not $\varphi$, using a notion of self-similarity that
  retains the intended physical meaning.

\item \textbf{Continuous zero-parameter physics.}
  Exhibit a theory with no adjustable dimensionless knobs whose
  state space is genuinely uncountable, and argue that
  continuity does not constitute hidden parameter freedom.

\item \textbf{Conservation without double-entry.}
  Produce a discrete conserved dynamics where conservation is
  representable without a balanced accounting structure.

\item \textbf{Pathological d'Alembert solution.}
  Construct a non-standard solution to~\eqref{eq:dalembert}
  satisfying the raw identity but violating regularity.
  This would not refute \cref{thm:t5} (regularity is explicit),
  but would show that regularity is a load-bearing hypothesis.
\end{enumerate}

% ===================================================================
\section{Conclusion}\label{sec:conclusion}
% ===================================================================

We have proved that any physics framework satisfying the five
admissibility conditions has a unique observational state on the
quotient: $\SQ(F) \simeq \mathbf{1}$.
The preferred scale is $\varphi = (1+\sqrt{5})/2$ and the cost
functional is $\Jcost(x) = \frac{1}{2}(x + x^{-1}) - 1$.

Three of the five assumptions are adopted axioms, not yet derived
from the zero-parameter posture alone.
The paper names six bridges (B1--B6) that would close the full
inevitability programme.
Of these, B2 (deriving the RCL/d'Alembert equation from ledger
structure) and B5 (constructing a unique prediction map) are
the deepest open problems.

The composition law is the d'Alembert functional equation
(\cref{sec:dalembert})---the unique symmetric composition law
on $(\Rp, \times)$ compatible with non-degeneracy.
It is not bespoke.
Whether it is \emph{forced} by the zero-parameter posture is
the central question the programme must answer.

% ===================================================================
\appendix
\section{Formal Verification}\label{app:lean}
% ===================================================================

All proved results are formally verified in Lean~4.
The following table maps paper results to Lean modules.
The Lean codebase uses one external axiom:
\texttt{aczel\_representation\_axiom} (the Acz\'el
classification, \cref{thm:aczel}).

\medskip
\begin{center}
\renewcommand{\arraystretch}{1.15}
\small
\begin{tabular}{lll}
\hline
\textbf{Paper result} & \textbf{Lean theorem} & \textbf{Module} \\
\hline
\cref{thm:phi} ($\varphi$ forcing) &
  \texttt{phi\_forced} &
  \texttt{PhiForcing} \\
\cref{thm:t5} (cost uniqueness) &
  \texttt{T5\_uniqueness\_complete} &
  \texttt{CostUniqueness} \\
\cref{thm:derived-uniform} (derived uniformity) &
  \texttt{zero\_params\_forces\_uniform} &
  \texttt{ModelIndependent} \\
\cref{thm:main} (quotient collapse) &
  \texttt{derived\_model\_independent\_exclusivity} &
  \texttt{ModelIndependent} \\
\cref{cor:equiv} (framework equivalence) &
  \texttt{model\_independent\_framework\_equiv} &
  \texttt{ModelIndependent} \\
B1 hierarchy theorem &
  \texttt{hierarchy\_forces\_fibonacci\_recurrence} &
  \texttt{HierarchyTheorem} \\
B4 ground-state corollary &
  \texttt{ground\_state\_implies\_cost\_minimum} &
  \texttt{GroundStateRestriction} \\
B6 type equivalence &
  \texttt{bridge\_B6\_observable\_type\_equiv\_trivial} &
  \texttt{ObservableTypeEquiv} \\
\hline
\end{tabular}
\end{center}

\medskip\noindent
All modules listed above compile with zero \texttt{sorry}.
The B2 scaffold (\texttt{RCLDerivation}) and B5 scaffold
(\texttt{PredictionMap}) contain explicit \texttt{sorry}-marked
open problems corresponding to \cref{op:rcl} and
\cref{sec:bridge-b5}.

% ===================================================================
\begin{thebibliography}{99}

\bibitem{Aczel1966}
J.~Acz\'el,
\emph{Lectures on Functional Equations and Their Applications},
Academic Press, 1966.

\bibitem{WZ2026}
J.~Washburn and M.~Zlatanovi\'c,
``Uniqueness of the Canonical Reciprocal Cost,''
\emph{Mathematics} \textbf{14}, 935 (2026).

\bibitem{Barghi2026}
J.~Washburn and A.~Rahnamai~Barghi,
``Reciprocal Convex Costs for Ratio Matching,''
\emph{Axioms} \textbf{15}(2), 151 (2026).

\bibitem{DAlembert2026}
J.~Washburn,
``D'Alembert Uniqueness is Fully Unconditional,''
Preprint (2026).

\end{thebibliography}

\end{document}
